% ============================================================
%  frontmatter/abstract-body.tex
%  Write your abstract here ONCE.
%  It is reused in both abstract.tex and form7a-abstract.tex
% ============================================================

Medical images acquired using X-ray and computed tomography (CT) are often
degraded by acquisition noise and patient motion, reducing diagnostic accuracy.
Deep learning methods have shown promise for image restoration, but existing
approaches have limitations. Standard denoisers often smooth out fine anatomical
details. Most activation functions apply the same operation everywhere, without
considering local context.

This thesis addresses these limitations through a three-phase investigation.
Phase~I presents \textbf{X-GAN}, a physics-informed generative adversarial
network for chest X-ray denoising, achieving a PSNR of 36.48~dB and an Edge
Preservation Index of 0.721 with only 2.17 million parameters. Phase~II
introduces \textbf{SAGA}, a spatially adaptive gated activation function
improving PSNR by 6.3~dB over ReLU on average across four architectures.
Phase~III will develop \textbf{Viveka}, a test-time refinement framework,
and \textbf{ADL$_1$}, a Sobolev-regularised loss function with an
attention-guided U-Net that guarantees gradient fidelity.